\startsection[title={Fermentação}]
	A fermentação ocupa posição de destaque na indústria biotecnológica, assumindo protagonismo, ainda hoje, em rotas produtivas de compostos essenciais para a sociedade \cite[authoryear][ref::metcalfe_0]. Grande parte das vacinas, medicamentos e biocombustíveis é obtida preferencialmente por processos fermentativos, e sua viabilidade econômica decorre justamente da rota empregada.

	Os microrganismos responsáveis pela fermentação podem ser bactérias, leveduras, arqueas, fungos ou células de mamíferos \cite[authoryear][ref::muller_0]. Existe uma ampla variedade de caminhos metabólicos possíveis considerando o conjunto desses organismos. As leveduras, por exemplo, são responsáveis pela produção de etanol, pão, bebidas alcoólicas, entre outros produtos, reforçando a relevância econômica e social desses microrganismos e de suas aplicações industriais.

	A produção de etanol por fermentação é um dos processos biotecnológicos mais importantes da indústria moderna, sendo amplamente utilizada para a obtenção de biocombustíveis e produtos químicos renováveis \cite[authoryear][ref::dashko_0]. Nesse processo, leveduras como \getbuffer[sacc] convertem açúcares provenientes de matérias-primas vegetais, como cana-de-açúcar, milho e biomassa lignocelulósica em etanol e dióxido de carbono sob condições anaeróbicas \cite[authoryear][ref::bhargav_0]. A fermentação alcoólica apresenta grande relevância econômica, especialmente em países como o Brasil, onde o etanol é utilizado em larga escala como combustível automotivo. Além de sua aplicação energética, o etanol também é empregado nas indústrias farmacêutica, alimentícia e química, muitas vezes na forma de gel ou líquida em diversas concentrações para ser utilizado na esterilização.

	\startplacefigure[title={Ilustração de células de microrganismos fermentativos}, reference={fig:cell_illustration}]
		\startcombination[nx=2]
			{\externalfigure[image/yeast_cell.png][height=200pt]}{(A) Levedura}
			{\externalfigure[image/bacterial_cell.png][height=200pt]}{(B) Bactéria}
		\stopcombination
	\stopplacefigure

	Na \in{Figura}[fig:cell_illustration] tem-se (A) uma célula de levedura e (B) a uma célula de bactéria. Nesta, vê-se algumas particularidades que a diferencia da célula de levedura, o que mais se destaca é a disposição do material genético (\chemical{DNA}), espalhado pelo citoplasma, caracterizando esses microrganismos como procariontes e os colocando em contraste com os microrganismos eucariontes de \chemical{DNA} concentrado no núcleo como a levedura. Existem alguns estudos como o de \cite[authoryears][ref::meadows_0] que demonstra o impacto da modificação genética nuclear de células eucariontes na sua performance. A mitocôndria também possui material genético, embora a extensão dos estudos relacionados à sua manipulação seja menor quando comparada à do núcleo. Um ponto importante na compreensão das reações químicas de oxirredução \cite[authoryear][ref::muller_0] é que elas ocorrem no citoplasma ou são coordenadas na mitocôndria quando associadas a processos aeróbicos. A diferenciação do interior de uma célula de microrganismos fermentativos e do exterior destes é importante para o tratamento de modelagens matemáticas dos seus metabolismos. Nestes modelos, as reações que acontecem no citoplasma ou na mitocôndria são balanceadas para que o fluxo molar total destes meios sejam nulos, possibilitando um tratamento homogêneo de sistema linear, e uma solução para os fluxos de todas reações intracelulares, de tal forma que os fluxos molares internos sejam retroalimentados em equações que modelam o balanço molar do exterior da célula, possivelmente por equações diferenciais \cite[authoryear][ref::sainz_0]. 
\stopsection

